% GENERATED FILE. Edit canonical lessons, then run npm run generate:books.
% canonical-source-hash: fnv1a64:21b0b5fda01921f4
% canonical-lessons: HI-C288-cori, HI-C288-curana, HI-C288-caukidar, HI-C288-gavah, HI-C288-bahar

\chapter{A Theft in the Lane}
\label{ch:hi-a-theft-in-the-lane}
\hlchaptermodality{288}

\begin{chapteropening}
\textbf{By the end of this chapter:} \emph{I can say what was stolen and who saw it.}

Complete the last lesson of chapter 288: I can say what was stolen and who saw it.
\end{chapteropening}

\section[\texorpdfstring{\hi{चोरी} (corī)}{चोरी (corī)}]{\hi{चोरी} (corī) — a theft}

\label{lesson:HI-C288-cori}



\begin{quote}
Before the new one: say the Hindi for turning over (in bed), then the Hindi for a stretch (of the body).
\end{quote}

\subsection*{You'll want to know: \hi{चोरी}}
\textbf{\hi{चोरी}} — \emph{corī} — \textquotedblleft{}a theft\textquotedblright{}.

\begin{quote}
\hi{चोरी} \hi{के} \hi{बाद} \hi{क्या} \hi{होगा}, \hi{यह} \hi{पुलिस} \hi{पर} \hi{निर्भर} \hi{करता} \hi{है।}
\end{quote}

\begin{quote}
\emph{corī ke bād kyā hogā, yah pulis par nirbhar kartā hai.}
\end{quote}

\begin{quote}
What happens after the theft depends on the police.
\end{quote}

Feminine.

\begin{cousinweb}
is inherited from Sanskrit \emph{caurikā}, \textquotedblleft{}theft\textquotedblright{} (CDIAL 4937). It sits beside \textbf{\hi{चोर}}, a thief: \textbf{\hi{चोर}} is the person, \textbf{\hi{चोरी}} is what the person does, and \textbf{\hi{चोरी} \hi{होना}} is for a theft to take place.
\end{cousinweb}

\begin{grammarlens}[title={What you've built}]
One more word for telling what happened.
\end{grammarlens}

\subsection*{Your turn}
\begin{itemize}
\raggedright
  \item \emph{Say it:} \emph{corī}
  \item \emph{Say it:} \emph{corī}, once more
  \item \emph{Say it:} \emph{corī ke bād kyā hogā, yah pulis par nirbhar kartā hai.} — What happens after the theft depends on the police.
  \item \emph{Recall:} say the Hindi for turning over (in bed), then the Hindi for a stretch (of the body), then say \emph{corī} again
\end{itemize}

\begin{tcolorbox}[breakable,colback=teal!4,colframe=teal!35!black,title={Before you move on}]
What does \hi{चोरी} mean? (\textquotedblleft{}A theft\textquotedblright{}.) Say it once more.
\end{tcolorbox}
\section[\texorpdfstring{\hi{चुराना} (curānā)}{चुराना (curānā)}]{\hi{चुराना} (curānā) — to steal}

\label{lesson:HI-C288-curana}



\begin{quote}
Before the new one: say the Hindi for a stretch (of the body), then the Hindi for a theft.
\end{quote}

\subsection*{You'll want to know: \hi{चुराना}}
\textbf{\hi{चुराना}} — \emph{curānā} — \textquotedblleft{}to steal\textquotedblright{}.

\begin{quote}
\hi{चोर} \hi{ने} \hi{मेरी} \hi{साइकिल} \hi{चुराई।}
\end{quote}

\begin{quote}
\emph{cor ne merī sāikil curāī.}
\end{quote}

\begin{quote}
The thief stole my bicycle.
\end{quote}

Past, with \hi{ने} (he or she …): \hi{उसने} \hi{बटुआ} \hi{चुराया}, \emph{usne baṭuā curāyā}.

\begin{cousinweb}
\textbf{\hi{चुराना}} (\emph{curānā}) belongs with Sanskrit \emph{corayati}, \textquotedblleft{}steals\textquotedblright{} (CDIAL 4933), where CDIAL records an older Hindi verb \emph{cornā}. With \textbf{\hi{ने}} the verb agrees with the thing taken: \textbf{\hi{साइकिल} \hi{चुराई}}, but \textbf{\hi{बटुआ} \hi{चुराया}}.
\end{cousinweb}

\begin{grammarlens}[title={What you've built}]
One more word for telling what happened.
\end{grammarlens}

\subsection*{Your turn}
\begin{itemize}
\raggedright
  \item \emph{Say it:} \emph{curānā}
  \item \emph{Say it:} \emph{curānā}, once more
  \item \emph{Say it:} \emph{cor ne merī sāikil curāī.} — The thief stole my bicycle.
  \item \emph{Recall:} say the Hindi for a stretch (of the body), then the Hindi for a theft, then say \emph{curānā} again
\end{itemize}

\begin{tcolorbox}[breakable,colback=teal!4,colframe=teal!35!black,title={Before you move on}]
What does \hi{चुराना} mean? (\textquotedblleft{}To steal\textquotedblright{}.) Say it once more.
\end{tcolorbox}
\section[\texorpdfstring{\hi{चौकीदार} (caukīdār)}{चौकीदार (caukīdār)}]{\hi{चौकीदार} (caukīdār) — a watchman}

\label{lesson:HI-C288-caukidar}



\begin{quote}
Before the new one: say the Hindi for a theft, then the Hindi for to steal.
\end{quote}

\subsection*{You'll want to know: \hi{चौकीदार}}
\textbf{\hi{चौकीदार}} — \emph{caukīdār} — \textquotedblleft{}a watchman\textquotedblright{}.

\begin{quote}
\hi{चौकीदार} \hi{रात} \hi{भर} \hi{गली} \hi{में} \hi{घूम} \hi{रहा} \hi{था।}
\end{quote}

\begin{quote}
\emph{caukīdār rāt bhar galī meṁ ghūm rahā thā.}
\end{quote}

\begin{quote}
The watchman was walking up and down the lane all night.
\end{quote}

Masculine.

\begin{cousinweb}
\textbf{\hi{चौकीदार}} (\emph{caukīdār}) joins \emph{caukī}, a guard post, from Sanskrit \emph{catuṣka}, \textquotedblleft{}a set of four, a square\textquotedblright{} (CDIAL 4629), with \emph{-dār}, usually traced to Persian \emph{dār}, \textquotedblleft{}holding\textquotedblright{}: the one who keeps the post. \textbf{\hi{दुकानदार}} is built with the same ending.
\end{cousinweb}

\begin{grammarlens}[title={What you've built}]
One more word for telling what happened.
\end{grammarlens}

\subsection*{Your turn}
\begin{itemize}
\raggedright
  \item \emph{Say it:} \emph{caukīdār}
  \item \emph{Say it:} \emph{caukīdār}, once more
  \item \emph{Say it:} \emph{caukīdār rāt bhar galī meṁ ghūm rahā thā.} — The watchman was walking up and down the lane all night.
  \item \emph{Recall:} say the Hindi for a theft, then the Hindi for to steal, then say \emph{caukīdār} again
\end{itemize}

\begin{tcolorbox}[breakable,colback=teal!4,colframe=teal!35!black,title={Before you move on}]
What does \hi{चौकीदार} mean? (\textquotedblleft{}A watchman\textquotedblright{}.) Say it once more.
\end{tcolorbox}
\section[\texorpdfstring{\hi{गवाह} (gavāh)}{गवाह (gavāh)}]{\hi{गवाह} (gavāh) — a witness}

\label{lesson:HI-C288-gavah}



\begin{quote}
Before the new one: say the Hindi for to steal, then the Hindi for a watchman.
\end{quote}

\subsection*{You'll want to know: \hi{गवाह}}
\textbf{\hi{गवाह}} — \emph{gavāh} — \textquotedblleft{}a witness\textquotedblright{}.

\begin{quote}
\hi{गवाह} \hi{ने} \hi{पुलिस} \hi{को} \hi{सब} \hi{कुछ} \hi{बताया।}
\end{quote}

\begin{quote}
\emph{gavāh ne pulis ko sab kuch batāyā.}
\end{quote}

\begin{quote}
The witness told the police everything.
\end{quote}

Masculine or feminine, by the person meant.

\begin{cousinweb}
\textbf{\hi{गवाह}} (\emph{gavāh}) is usually traced to Persian \emph{gavāh}, \textquotedblleft{}a witness\textquotedblright{}, a word that reached Hindi through the courts. It is used of a man or a woman: \textbf{\hi{वह} \hi{गवाह} \hi{है}}.
\end{cousinweb}

\begin{grammarlens}[title={What you've built}]
One more word for telling what happened.
\end{grammarlens}

\subsection*{Your turn}
\begin{itemize}
\raggedright
  \item \emph{Say it:} \emph{gavāh}
  \item \emph{Say it:} \emph{gavāh}, once more
  \item \emph{Say it:} \emph{gavāh ne pulis ko sab kuch batāyā.} — The witness told the police everything.
  \item \emph{Recall:} say the Hindi for to steal, then the Hindi for a watchman, then say \emph{gavāh} again
\end{itemize}

\begin{tcolorbox}[breakable,colback=teal!4,colframe=teal!35!black,title={Before you move on}]
What does \hi{गवाह} mean? (\textquotedblleft{}A witness\textquotedblright{}.) Say it once more.
\end{tcolorbox}
\section[\texorpdfstring{\hi{बाहर} (bāhar)}{बाहर (bāhar)}]{\hi{बाहर} (bāhar) — outside}

\label{lesson:HI-C288-bahar}



\begin{quote}
Before the new one: say the Hindi for a watchman, then the Hindi for a witness.
\end{quote}

\subsection*{You'll want to know: \hi{बाहर}}
\textbf{\hi{बाहर}} — \emph{bāhar} — \textquotedblleft{}outside\textquotedblright{}.

\begin{quote}
\hi{चोरी} \hi{के} \hi{बाद} \hi{हम} \hi{सब} \hi{घर} \hi{से} \hi{बाहर} \hi{आए।}
\end{quote}

\begin{quote}
\emph{corī ke bād ham sab ghar se bāhar āe.}
\end{quote}

\begin{quote}
After the theft we all came out of the house.
\end{quote}

\begin{cousinweb}
\textbf{\hi{बाहर}} (\emph{bāhar}) is inherited from \emph{*bāhira}, \textquotedblleft{}outer\textquotedblright{}, a form CDIAL reconstructs from its descendants (CDIAL 9226). It pairs with \textbf{\hi{अंदर}}: \textbf{\hi{घर} \hi{से} \hi{बाहर}}, out of the house, against \textbf{\hi{घर} \hi{के} \hi{अंदर}}, inside it.
\end{cousinweb}

\begin{grammarlens}[title={What you've built}]
One more word for telling what happened.
\end{grammarlens}

\subsection*{Your turn}
\begin{itemize}
\raggedright
  \item \emph{Say it:} \emph{bāhar}
  \item \emph{Say it:} \emph{bāhar}, once more
  \item \emph{Say it:} \emph{corī ke bād ham sab ghar se bāhar āe.} — After the theft we all came out of the house.
  \item \emph{Recall:} say the Hindi for a watchman, then the Hindi for a witness, then say \emph{bāhar} again
\end{itemize}

\begin{tcolorbox}[breakable,colback=teal!4,colframe=teal!35!black,title={Before you move on}]
What does \hi{बाहर} mean? (\textquotedblleft{}Outside\textquotedblright{}.) Say it once more.
\end{tcolorbox}
